% ===================================================================
%  第 4 节  权重零处的月光（Moonshine at weight zero）
%  原文：source/original/TASI10.tex 第 1482--1762 行（PDF p.38--45）
%  说明：公式环境连同 \label 逐字照录；行内公式照抄；散文与标题译为中文。
% ===================================================================
\section{权重零处的月光}
\label{cha:weight0}

\begin{keytakeaways}
  \item 权重零的月光有两个实例：魔群月光联系魔群 $\MM$ 的表示与模函数 $J(\tau)=q^{-1}+196884q+\cdots$，Conway 月光则联系 $\Co_0$ 与某些亏格零子群的主模（Hauptmodul）。
  \item 魔群月光的核心对象是 McKay–Thompson 级数 $T_g(\tau)=\sum_{n\geq-1}\tr_{V_n}(g)q^n$；猜想断言每个 $T_g$ 都等于展开为 $q^{-1}+O(q)$ 的唯一的主模。
  \item 魔群模 $V^\natural$ 由 Leech 格 VOA 作 $\ZZ_2$ 轨道流形得到：$c=24$、无重量一态，$196884$ 维的重量二空间正是 Griess 代数，且 $\Aut(V^\natural)\cong\MM$。
  \item 可复制性公式 \eq{eq:replicability_J} 给出主模系数的递推；Borcherds 把它解释为魔群 Lie 代数 $\mathfrak{m}$ 的分母恒等式，由此证明魔群月光猜想。
  \item Conway 侧对应的命题是定理 \ref{thm:conway_moonshine}：$T^s_g(2\tau)$ 是亏格零主模；扭模上的迹 $T^s_{g,\text{tw}}$ 在 $g$ 于 $\Lambda$ 上有不动点时为常数 $-\chi_g$，否则也是主模。
\end{keytakeaways}

本节回顾权重为零的两种月光联系：魔群月光与 Conway 月光。就模对象的指定、对称性的来源以及物理背景而言，它们是当前理解得最清楚的月光实例。

% [Tutor Note: “权重零”指相关模对象是权 $0$ 的模函数，更准确地说，是亏格零子群的主模：其 $q$ 展开只有 $q^{-1}$ 极点而没有更高阶极点。正因如此，一旦亏格零群 $\Gamma_g$ 确定，函数 $T_g$ 就被唯一决定；这是权零情形“最清楚”的确切含义。]

\subsection{魔群月光}
\label{sec:mm}

魔群月光可以说是上世纪数学中最迷人的篇章之一：有限群与模对象正是最早在这里被发现通过物理结构相互关联。随着月光理论的发展，我们相信，从多种角度看，魔群月光仍将是月光现象中最杰出的范例。本节简要描述魔群月光的特征；关于这段优美故事的其它优秀综述，特别是其历史侧面，参见 \cite{MR2201600,mason2014vertex,MR3375653} 及其中的参考文献。

% [Tutor Note: 作者说“通过物理结构（via physical structures）相互关联”，指的是以魔群为对称群的二维手征共形场论：魔群模 $V^\natural$ 既是 VOA，又是某个手征 CFT 的态空间，$c=24$。这一物理构造迟至 1988 年才由 Frenkel–Lepowsky–Meurman 给出，正是本节后半“月光模”一段的主题。]

\emph{魔群月光}（monstrous moonshine）一词由 \cite{conway_norton} 创造，指魔群 $\MM$ 的表示论与模形式
\begin{equation}\label{Jcoeff}
J(\tau)=\sum_{n\geq-1}a_n \, q^n=q^{-1}+196884q+21493760q^2+864299970q^3+\cdots~,
\end{equation}
之间的意外联系；我们在 \eq{eq:def_J} 中已经遇到过这个模形式。魔群月光的发展始于 McKay 的关键观察：$J$ 的 $q$ 展开中系数 $196884$ 可分解为 $196884=1+196883$，其中两个加数正是 $\MM$ 最小的两个不可约表示的维数。Thompson 在 \cite{Tho_NmrlgyMonsEllModFn} 中对接下来几个系数观察到类似的分解：
\begin{equation}
\label{eq:J_decomp}
\begin{split}
1&=1 \\
196884&=196883+1 \\
21493760&=21296876+196883+1 \\
864299970&=842609326+21296876+2\cdot196883+2\cdot1 \\
&\cdots 
\end{split}
\end{equation}
其中 $1$、196883、21296876 与 842609326 是 $\MM$ 的某些不可约表示的维数。这一观察引出了如下猜想：存在无限维的 $\ZZ$ 分次魔群模
\begin{equation}
V=\bigoplus_{n\geq-1}V_n~,
\end{equation}
使得对一切 $n\geq-1$ 都有 $\dim V_n=a_n$。换言之，$J$ 获得了 $V$ 的分次维数这一解释
\begin{equation}
\label{eq:J_graded_trace}
J(\tau)=\sum_{n\geq-1}\dim V_n q^n.
\end{equation}
注意 $V_0$ 是空的，这对应于 $J$ 的常数项为零。

% [Tutor Note: $196884=196883+1$ 不是数值巧合：$V^\natural$ 的重量二空间（$196884$ 维）正是 Griess 代数，一个交换而非结合的代数；魔群在其 $196883$ 维子空间上忠实作用，多出的 $1$ 维由 Virasoro 向量（$L_{-2}|0\rangle$ 型态）张成且为魔群不变。所以 $196884-1=196883$ 直接反映了“Virasoro 单态加上魔群的忠实作用”这一结构。]

% [Tutor Note: “$V_0$ 是空的”指 $a_0=0$，即不存在 $n=0$ 的态。这里的分次指标是 $n=L_0-c/24$（魔群模 $c=24$）而不是 $L_0$：真空落在 $n=-1$（给出 $q^{-1}$），重量一（$L_0=1$）的态本该落在 $n=0$。故 $V_0$ 为空恰好就是“无重量一态”，对应文中后面对 $V^\natural$ 的讨论。]

\myparagraph{猜想。}
如上陈述的猜想并不有趣：既然所有 $a_n$ 都是非负整数，只要让每个 $V_n$ 含 $c(n)$ 份 $\MM$ 的平凡表示，就可使 \eq{eq:J_graded_trace} 成立。为了获取 $\MM$-作用的信息，Thompson 在 \cite{Tho_FinGpsModFns} 中还提议考察 $V$ 的分次特征标，即所谓的 \emph{McKay–Thompson 级数}（McKay-Thompson series），其定义为
\begin{equation}
T_g(\tau)\coloneqq\sum_{n\geq-1}\tr_{V_n}(g)q^n~,
\end{equation}
对每个元素 $g\in\mathbb M$ 定义（并有 $T_e=J$）。注意 $q$ 级数 $T_g(\tau)$ 的常数项也必须为零。从定义显然可见，$T_g$ 是类函数，即 $T_g=T_{hgh^{-1}}$。因此最多有 $194$ 个不同的 $T_g$，因为 $\MM$ 有 $194$ 个共轭类。事实上，$T_g$ 只给出 $171$ 个不同的函数。魔群月光的要点在于：这些分次迹函数也表现出模性质，而且是唯一一类没有常数项的主模（参见 \S\ref{subsec:modular_forms}），正如 Conway 与 Norton 在 \cite{conway_norton} 中提出的下述惊人猜想所述：

% [Tutor Note: 原文此处用 $c(n)$ 表示平凡表示的重数，但前文只定义了 $a_n$；按上下文 $c(n)$ 应即 $a_n$。作者的意思是：分次维数只记录重数，任何非负整数都可以充当重数，所以它完全不含 $\MM$-作用的信息——这正是必须考察分次特征标 $T_g$ 的理由。]

% [Tutor Note: $T_g$ 的定义有三处要点：取 $V_n$ 上的迹而非维数（这才记录群作用）、级数按 $q^n$ 展开（$n$ 即 $L_0-c/24$，真空已在 $n=-1$）、常数项必须为零。所谓“从定义显然可见”的依据只是迹的共轭不变性，即 $\tr_{V_n}(hgh^{-1})=\tr_{V_n}(g)$。$194$ 个共轭类只给出 $171$ 个不同函数：除 $g$ 与 $g^{-1}$ 因迹取整数值而给出同一函数外，还有若干对共轭类发生更强的重合。]

{\conj \label{conj:conwa_norton}  {\bf（魔群月光猜想）}\\
对每个 $g\in \MM$，McKay–Thompson 级数 $T_g$ 与唯一的主模 $J_{\Gamma_g}$ 重合，后者在 $\t\to i\infty$ 附近有展开 $q^{-1}+O(q)$，其中 $\Gamma_g\leq\SL_2(\RR)$ 是某个亏格零子群。此外，每个 $\Gamma_g$ 都以 $\Gamma_0(N)$ 为正规子群，其中 $N$ 整除量 $|g|\gcd(24,|g|)$。
}

% [Tutor Note: 猜想的关键词是“唯一”：亏格零群 $\Gamma_g$ 给定之后，满足 $q^{-1}+O(q)$ 展开的主模在相差常数的意义下唯一，而“无常数项”恰好把这个常数固定下来。故猜想断言的不止是 $T_g$ 的模性，而是 $T_g$ 被 $\Gamma_g$ 完全决定。]

~\\~\\
鉴于这个猜想的重要性，我们暂停下来作几点评论。注意 $\Gamma_g$ 常常不是 $\SL_2(\ZZ)$ 的子群；只有对某些 $g$ 才有 $\Gamma_g= \Gamma_0(N)$（参见 \eq{def:cong_subgr}），其中 $N$ 满足上述条件。一般地，$\Gamma_g$ 是 $\Gamma_0(N)$ 在 $\SL_2(\RR)$ 中的正规化子，其中通常涉及所谓的 Atkin–Lehner 对合。为后文计，我们会特别关心如下形式的群
\be\label{eqn:groupAL}
\Gamma^{N+K} :=
	\left\{
		\frac1{\sqrt{n}}\begin{pmatrix}
		an&b\\cN&dn
		\end{pmatrix}
		\mid
		adn-bcN/n=1,\;
		 n \in K
	\right\},
\ee
其中 $K<{\rm Ex}_N$ 是 $N$ 的精确除子（exact divisor）群的子群。称 $e$ 是 $N$ 的精确除子，若 $e|N$ 且 $(f,{N\over f})=1$；这些除子按乘法 $f\ast f'  = {ff' \over (f,f')^2}$ 成群。一个特别简单的情形是 $N$ 为素数 $p$，此时完全正规化子（对应 $K=\{1,p\}$）由下式给出
\be
\Gamma_0(p)+:= \Big\langle \Gamma_0(p), {1\over \sqrt{p}} \bem 0&-1\\ p& 0\eem \Big\rangle ~. 
\ee

% [Tutor Note: 原文此处写 $(f,{N\over f})=1$，但被定义的对象是 $e$；按上下文应为 $(e,{N\over e})=1$（PDF p.39 同此写法）。疑为笔误，译文照录原式，读者应按 $(e,{N\over e})=1$ 理解。]


魔群月光的一个先兆，早于 McKay 的观察，是 Ogg 的如下观察 \cite{ogg1974automorphismes}。他指出，$\Gamma_0(p)+$ 在上半平面上定义亏格零商，当且仅当
\be
p\in \{ 2, 3, 5, 7, 11, 13, 17, 19, 23, 29, 31, 41, 47, 59, 71\}, 
\ee
而这恰好是整除魔群阶的素数集合；随后他向任何能解释这一巧合的人悬赏一瓶 Jack Daniel's 威士忌 \cite{ogg1974automorphismes}。魔群月光之所以能照亮这一神秘巧合，是因为所有亏格零 $\Gamma_0(p)+$ 的主模都在月光中作为某个阶为 $p$ 的 $g\in \MM$ 的 McKay–Thompson 级数 $T_g$ 出现。当 $\Gamma_g\not\subset \SL_2(\ZZ)$ 时，CFT 的模性无助于解释 $\Gamma_g$ 的模性何以出现，因为在 CFT 中模性来自环面的映射类群（参见 \S\ref{sec:CFT}）。

% [Tutor Note: 这 $15$ 个素数恰是整除 $|\MM|$ 的全部素因子，也恰是超奇异素数（supersingular primes）：对素数 $p$，$\Gamma_0(p)+$ 亏格零当且仅当 $p$ 为超奇异素数。这正是 Ogg 观察到的巧合，后来由魔群月光给出解释；注意列表止于 $71$，更大的素数都不满足条件。]

魔群月光的关键亏格零性质获得了一种有用的改写 \cite{DunFre_RSMG}：这些主模可以（在相差一个常数的意义下）作为 \emph{Rademacher 和}（Rademacher sum）得到，即把极点项（此处为 $q^{-1}$）在 $\SL_2(\RR)$ 的适当子群（此处为 $\Gamma_g$）作用下的像作正则化求和。这一 \emph{Rademacher 可和性}（Rademacher summability）随后在本影月光的发现中起了关键作用（参见 \S\ref{sec:umbral}）。最近，亏格零性质得到了解释：在魔群理论的弦论实现中，$\Gamma_g$ 扮演弦论对称群的角色，而对理论的超对称指标要求一个物理的解析性条件。详见 \S\ref{sec:physics_monstrous}。


\myparagraph{月光模。}
这个猜想由 Atkin、Fong 与 Smith 用数值方法验证（参见 \cite{fong1980characters,smith1985head}），其思路来自 Thompson（参考文献见 \cite{MR2201600}）。更具体地说，他们证明：对每个 $n\ge -1$，\cite{conway_norton} 中指定的各主模的 $q^n$ 系数，都与 $\MM$ 的某个虚拟表示的特征标重合。后来 Frenkel、Lepowsky 与 Meurman \cite{FLMPNAS,FLMBerk} 给出了构造性的验证：他们显式构造了一个魔群模 $V=V^\natural$。该模具有全纯 VOA 的结构：它只有一个不可约模，就是它自身。


构造 $V^\natural$ 的出发点是 $24$ 个手征玻色子 $X^i(z)$，紧化在由 Leech 格 $\Lambda$ 定义的 $24$ 维环面 $\RR^{24}/\Lambda$ 上。这得到一个中心荷 $c=24$ 的 VOA $V(\Lambda)$，其配分函数的 $q$ 展开以
\( Z_{V(\Lambda)}(\tau)=q^{-1}+\dots\) 开始。这一点连同模不变性，把该函数确定为与 $J(\tau)$ 相差一个加性常数。同时我们知道这个常数必须是多少，因为 Leech 格没有根向量，故 $\Theta_{\Lambda}(\t) =\sum_{v\in \Lambda} q^{\langle v,v\rangle/2}  = 1+ O(q^2)$，从而
\be
Z_{V(\Lambda)}(\tau)= {\Theta_{\Lambda}(\t)\over \eta^{24}(\t)} = J(\t)+24~. 
\ee
换言之，由于 Leech 格无根，形如 $e^{ik\cdot\Phi}$ 的格顶点算子的权都大于一，剩下的唯一权一原初场就是 $24$ 个场 $\partial X^i$。

% [Tutor Note: 常数 $24$ 极易漏掉：$\Theta_\Lambda=1+O(q^2)$ 给出 $q^0$ 系数 $1$，而 $1/\eta^{24}=q^{-1}(1+24q+\cdots)$，故 $\Theta_\Lambda/\eta^{24}=q^{-1}+24+\cdots$，正是 $J+24$。若直接写 $Z=J$ 就丢掉这 $24$，模性质也随之破坏——只有 $\Theta_\Lambda/\eta^{24}$ 在 $\SL_2(\ZZ)$ 下不变。]

为了与 $J$ 精确匹配，我们希望去掉这些原初场，这可以通过对 $V(\Lambda)$ 作一个简单的 $\ZZ_2$ 轨道流形实现：作用为 $X^i\rightarrow-X^i$，对应 $\Lambda$ 的 $\lbrace{\rm id},-{\rm id}\rbrace\cong\ZZ_2$ 对称性，它包含在 $\Aut(\L)\cong\Co_0$ 中。确实，我们可以显式地算出轨道流形理论的配分函数如下。注意 $24$ 个手征玻色子的 $\ZZ_2$-twined 配分函数为
\be
Z\left(\mbox{\fontsize{9}{10}\selectfont $-$}\underset{+}{\Box}~;\t\right)= {1\over q\prod_{n>0}(1+q^{n})^{24}} = \left(2 \eta(\t)\over \theta_2(\t) \right)^{12}. 
\ee
轨道流形意味着 $V^\natural$ 是未扭结扇区与扭结扇区的 $\ZZ_2$-不变投影的直和（参见 \eq{eqn:orb_PF}）。由 \eq{eqn:modular_boundary_condition} 与 \eq{eqn:transf_theta} 可得
\begin{gather}
\begin{split}
Z_{V^\natural}(\tau)& ={1\over 2}\left( J(\t)+24 + \left(2 \eta(\t)\over \theta_2(\t) \right)^{12} -  \left(2 \eta(\t)\over \theta_3(\t) \right)^{12} +\left(2 \eta(\t)\over \theta_4(\t) \right)^{12} \right) \\& = J(\t). 
\end{split}
\end{gather}

% [Tutor Note: 四个 $\theta$ 项前的符号依次为 $+,-,+$，来自 $\ZZ_2$ 轨道流形的四种边界条件（未扭结、twined、扭结、扭结-twined）；$\theta_2,\theta_3,\theta_4$ 分别对应 $\tau$ 平移下的不同自旋结构。漏掉任一符号、或把 $\theta_2$ 与 $\theta_3$ 的角色弄反，都会使 $V^\natural$ 的配分函数不再等于 $J$，从而破坏模不变性。]

还需说明 $\Aut(V^\natural)$ 就是魔群。注意 $\Aut V(\Lambda)$ 有一个连续部分 $T$，它是 $24$ 维环面，对应手征玻色子的平移对称性以及 $24$ 个权一原初场 $\pa X^i$。总的对称性由（非分裂的）短正合列
\be
1 \rightarrow T \rightarrow {\rm Aut}V(\Lambda) \rightarrow \Co_0 \rightarrow 1~. 
\ee 
% [Tutor Note: 原文此处写作 \rm{Aut}V(\Lambda)（老式字体命令），在 XeTeX+ctex 下会使紧随其后的 \Lambda 丢失字形（编译日志出现 Missing character U+0003，PDF 中显示为 V()）；译文改为 {\rm Aut}V(\Lambda)，数学内容与原文完全相同，仅修正该排版缺陷。]
刻画。$\ZZ_2$-轨道流形把自同构群破缺为其离散部分 $2^{24}.\Co_0$，后者保持分解 $V^\natural=V^\natural_+\oplus V^\natural_-$，并且与 $\MM$ 的某个极大子群 $2^{1+24}.\Co_1$ 有提示性的相似。从 $V^\natural_+$ 与 $V^\natural_-$ 对 $V^\natural$ 中权二（以及类似地权三、权四、$\dots$）态的贡献可以清楚看出，魔群必须混合二者，因而不能分别保持（未）扭结扇区。注意 $V^\natural$ 的 $196884$ 维权二态空间具有交换而非结合代数的结构（任何 VOA 都如此），并且可以证明它正是 1980 年构造的、曾用来构造魔群本身的 \emph{Griess 代数}（Griess algebra）\cite{Gri_FG}。由此以及 $V^\natural$ 的 VOA 结构可以证明 $\Aut(V^\natural)$ 确实是魔群，并且可以通过把某个混合 $V^\natural_\pm$ 的二阶对称性添加进 $\Aut(V(\Lambda))$ 的离散部分而得到。


\myparagraph{魔群月光的证明。}
为证明 Frenkel、Lepowsky 与 Meurman 构造的 $V^\natural$ 确实“胜任”，需要证明
\be
T_g^{V^\natural}(\tau):= \tr_{V^\natural} g \, q^{L_0-c/24} 
\ee
与 \cite{conway_norton} 中指定的相应主模 $J_{\Gamma_g}$ 重合。已知主模的系数满足某些递推公式，只需少数几个系数就能确定全部系数。在最简单的情形，这些递推公式由一个著名的恒等式编码：
\begin{equation}
\label{eq:replicability_J}
p^{-1}\prod_{m>0\atop n\in\ZZ}\left(1-p^mq^n\right)^{a_{mn}}=J(\rho)-J(\tau)~,
\end{equation}
它由 Zagier、Borcherds 等人独立发现。这里 $p=e^{2\pi i\rho}$，而 $a_i$ 表示 $J$ 的 $q$ 展开中的 $q^i$ 系数（参见（\ref{Jcoeff}））。这个恒等式给出 $a_i$ 之间无穷多个关系，使人们能仅由 $a_1,a_2,a_3,a_5$ 完全确定所有系数。显然，只要能证明 $T_g^{V^\natural}(\tau)$ 的系数之间存在同类型恒等式，并显式比较确定两侧整个函数所需的少数几个系数，证明即可完成。

% [Tutor Note: 可复制性公式 \eq{eq:replicability_J} 的指标约定最易看错：$a_{mn}$ 的下标是乘积 $mn$ 这一个整数，不是两个独立指标，即 $a_{mn}$ 就是 $J$ 展开中 $q^{mn}$ 的系数。乘积取遍 $m>0$、$n\in\ZZ$，$n<0$ 的因子也参与并负责极点结构。由于 $a_i$ 可为负，$(1-p^mq^n)^{a_{mn}}$ 须按形式幂级数的二项式展开理解。]

% [Tutor Note: 作者略去了“为何 $a_1,a_2,a_3,a_5$ 四个系数足够”的论证：把 \eq{eq:replicability_J} 按 $p^mq^n$ 逐级比较系数，会得到关于 $a_N$（$N=mn$）的递推关系，右端只在与 $J(\rho)$ 或 $J(\tau)$ 的单项式对齐时非零；于是给定上述初始值即可递归地确定其余全部系数。具体递推与收敛性见 Borcherds 1992 的原文及 Zagier 的综述。]

这正是 Borcherds 所做的。他通过引入 \emph{广义 Kac–Moody 代数}（generalised Kac-Moody algebra）的概念得到了可复制性公式；这类代数可看成允许虚单根的 Kac–Moody 代数的推广。随后他构造了一个广义 Kac–Moody 代数 $\mathfrak{m}$（也称“Borcherds–Kac–Moody 代数”），即所谓的魔群 Lie 代数（Monster Lie algebra）。粗略地说，这一构造通过研究一个类 BRST 算子的上同调完成，该算子作用在 $V^\natural\times\Gamma^{1,1}$ 上，其中 $\Gamma^{1,1}$ 是符号为 $(1,1)$ 的唯一幺模格。这一构造在弦论中有自然的解释：考虑以 $V^\natural$ 为背景的二次量子化弦（second quantised strings）。详见 \S\ref{sec:physics_monstrous}。

% [Tutor Note: Borcherds 的证明路线可概括为：广义 Kac–Moody 代数允许虚单根，其 Weyl 分母因此是无穷乘积；把 Weyl–Kac 特征标公式用于平凡表示，即得分母恒等式。技术上先在 $26$ 维偶幺模格 $II_{25,1}$（Leech 格与 $\Gamma^{1,1}$ 的直和）上构造“伪魔群 Lie 代数”（Fake Monster Lie algebra）以检验这一机制，再把 $24$ 维部分换成 $V^\natural$，得到魔群 Lie 代数 $\mathfrak{m}$，其根重数恰为 $J$ 的系数 $a_{mn}$。]


Borcherds 成功地把可复制性公式 \eqref{eq:replicability_J} 导为他附在 $V^\natural$ 上的魔群 Lie 代数 $\mathfrak{m}$ 的\emph{分母恒等式}。与通常的 Kac–Moody 代数一样，分母恒等式来自把 Lie 代数的 Weyl–Kac 特征标公式用于平凡表示；在现在的情形，它把无穷和与无穷乘积联系起来，正是我们在 \eq{eq:replicability_J} 中看到的结构。此外，通过考察 $V^\natural$ 上的 $\MM$-作用，还可以由 $\mathfrak{m}$ 得到类似恒等式
\begin{equation}
\label{eq:replicability_T}
p^{-1}\exp\left[-\sum_{k>0}\sum_{m>0\atop n\in\ZZ} a_{mn}^{g^k}\frac{p^{mk}q^{nk}}{k}\right]=J_{\Gamma_g}(z)-J_{\Gamma_g}(\tau)~,
\end{equation}
为其它主模所满足，其中 $a_i^g$ 是 $J_{\Gamma_g}$ 的 $q$ 展开系数。综合以上各部分，便证明了 $V^\natural$ 作为魔群月光模的正当性。

% [Tutor Note: 式 \eq{eq:replicability_T} 中的上标 $g^k$ 来自对 $g$ 的各次幂取迹：可复制性要求同时掌握 $T_{g^k}$ 的系数，故乘积的每个因子由 $a^{g^k}$ 控制并带 $1/k$ 权重——这正是“可复制性”一词的来源。注意 $k$ 同时出现在 $p^{mk}q^{nk}$ 的指数中，不能与 $m,n$ 混淆；对每个固定次数的项只有有限多个 $k$ 贡献，故该式是良定义的形式幂级数恒等式。]

\myparagraph{广义魔群月光。}
在 \cite{generalized_moonshine} 中，Norton 以\emph{广义魔群月光}（generalised monstrous moonshine）之名提出了魔群月光的一个推广。他设想存在一条规则：对每个元素 $g\in\MM$ 指定中心化子群 $C_\MM(g)$ 的一个分次射影表示 $V(g)=\bigoplus_{n\in\QQ}V(g)_n$，并对每一对交换元素 $(g,h)$ 指定上半平面 $\HH$ 上的一个全纯函数 $T_{(g,h)}$，满足以下条件：
\begin{enumerate}[label=(\roman*)]
\item 
$~T_{\left(g^ah^c,g^bh^d\right)}(\tau)=\gamma ~T_{(g,h)}\left(\frac{a\tau+b}{c\tau+d}\right)~~~~$ 其中 $~~\begin{pmatrix} a &b \\ c&d\end{pmatrix}\in\SL_2(\ZZ)$，而 $\gamma$ 为 $24$ 次单位根。
\item $~T_{\left(g,h\right)}(\tau)=T_{\left(k^{-1}gk,k^{-1}hk\right)}(\tau)~~~~$ 其中 $k\in\MM$。
\item 存在 $h$ 在 $V(g)$ 上的一个提升 $\tilde{h}$，作为线性变换，使得
\begin{equation}
T_{(g,h)}(\tau)=\sum_{n\in\QQ}\tr_{V(g)_n}\left(\tilde{h}~q^{n-1}\right)~.
\end{equation}
\item \label{haupt} $T_{(g,h)}(\tau)$ 或为常数，或为 $\SL_2(\ZZ)$ 的某个亏格零同余子群的主模。
\item $T_{(e,h)}(\tau)$ 与 $T_h(\tau)$ 重合，后者是魔群月光附给 $h\in\MM$ 的 McKay–Thompson 级数。
\end{enumerate}


如 \S\ref{sec:orbifolds} 的讨论所示，除 \ref{haupt} 之外，上述所有性质都可以在应用于 $V^\natural$ 的全纯轨道流形（holomorphic orbifolds）\cite{DixonL.1988Batb} 框架中理解。特别地，函数 $T_{(g,h)}$ 可以看作扭模 $V^\natural_g$ 的 $h$-twined 特征标。广义魔群月光的证明最近由 \cite{carnahan2012generalized} 完成，其中对每个 $g\in \MM$ 都构造了一个广义 Kac–Moody 代数 ${\mathfrak m}_g$，它是魔群 Lie 代数 $\mathfrak{m}$ 的推广。


\subsection{Conway 月光}
\label{sec:conway_moonshine}

Conway 月光建立了 $\Co_0$ 与 $\SL_2(\RR)$ 的某些亏格零子群的主模之间的联系；$\Co_0$ 通过 $\Co_1\cong\text{Co}_0/\lbrace\pm\text{Id}\rbrace$ 与 Conway 的散在单群 $\Co_1$ 相联系。

回忆 \S\ref{sec:sporadic_lattices}：$\Co_0$ 同构于 Leech 格 $\Lambda$ 的自同构群。在此背景下，Conway 月光的迹象已经出现在最初的魔群月光论文 \cite{conway_norton} 中：作者为元素 $g\in\text{Co}_0$ 指定了亏格零群 $\Gamma_g<\SL_2(\mathbb{R})$。设 $\left\lbrace\lambda_i,\lambda_i^{-1}\right\rbrace_{i=1}^{12}$ 为 $g$ 在 Leech 格 $\Lambda\otimes_{\mathbb{Z}}\mathbb{C}$（嵌入某个复向量空间）上自然作用的 $24$ 个本征值，则 $\Gamma_g$ 由全纯函数
\begin{equation}
t_g(\t):=q^{-1}\prod_{n>0}\prod_{i=1}^{12}\left(1-\lambda_iq^{2n-1}\right)\left(1-\lambda_i^{-1}q^{2n-1}\right)=q^{-1}-\chi_g+\mathcal{O}(q)~. 
\end{equation}
的不变群给出。注意 $\chi_g = \sum_i (\lambda_i+\lambda_i^{-1})$ 一般不为零，故与上一小节讨论的魔群月光函数不同，$t_g$ 的常数项非零。


Conway 月光则引入一个 Conway 模 $V^{s\natural}$，其 McKay–Thompson 级数与常数项为零的主模重合。它由 \cite{DuncanJohnF.2007SfCl, DuncanJohnf.r.2015TMMF} 发展起来（关于该构造的精彩综述另见 \cite{MR3465528}、\cite{Cheng:2014owa} 与 \cite{Harrison:2018joy}）。Conway 模 $V^{s\natural}$ 在同构意义下是唯一的超顶点算子代数（super-VOA，SVOA），具有 $c_{V^{s\natural}}=12$ 与 $\mathcal{N}=1$ 超共形结构，且不含重量 $1/2$ 的态。它可以构造为 $E_8$ 环面上八个玻色子连同其费米超伙伴的理论的 $\ZZ_2$ 轨道流形；也可以构造为 $24$ 个自由手征费米子 $k_a$（$a=1,2,\dots,24$）在 $k_a \to -k_a$ 作用下的 $\ZZ_2$ 轨道流形。这可以与魔群月光模 $V^\natural$ 相比：相应的魔群模 $V^\natural$ 由 Leech 格 VOA（$\RR^{24}/\Lambda$ 上紧化的 $24$ 个手征玻色子）的 $\ZZ_2$ 轨道流形构造，所得理论有 $c_{V^\natural}=24$ 且无重量 $1$ 的态。事实证明，$V^{s\natural}$ 与 K3 曲面的弦论对称性有有趣的联系（参见 \S\ref{subsec:conwayK3}）。下面我们给出 Conway 模 $V^{s\natural}$ 的更多细节。


考虑 $24$ 个实手征费米子 $k_a$ 及相应的复费米子
\begin{equation}
\psi^\pm_j=\frac{1}{\sqrt{2}}\left(k_{2j-1}\pm ik_{2j}\right)~,~~~j=1,\dots,12~, 
\end{equation}
其非零 OPE 与应力张量为
\begin{equation}
\psi^\pm_i(z)\psi^\mp_j(w)\sim\frac{\delta_{ij}}{z-w}~~,~~~~~L=-\frac{1}{2}\sum_{i=1}^{12}\normord{\psi_i^+\partial\psi_i^- + \psi_i^-\partial\psi_i^+}~.
\end{equation}


记 $\mathfrak{a}$ 为由这些费米子张成的 $24$ 维向量空间。由于费米子允许周期与反周期边界条件，理论中存在两个扇区。反周期（Neveu–Schwartz）扇区只含一个基态 $|0\rangle$ 以及半整数权重的激发；周期（Ramond）扇区含整数权重的激发，并有 $2^{12}$ 个简并基态。这种简并来自零模所满足的 Clifford 代数
\begin{equation}
\left\lbrace\psi^\pm_{i,0},\psi^\pm_{j,0}\right\rbrace=0~,~~~\left\lbrace\psi^\pm_{i,0},\psi^\mp_{j,0}\right\rbrace=\delta_{ij}~,
\end{equation}
而这些零模还与 $L_0$ 对易。于是可以在满足 $\psi^+_{i,0}|s\rangle=0$ 的基态 $|s\rangle$ 上作用 $\psi^-_{i,0}$，构造出 Ramond 基态。也就是说，Ramond 基态由单项式
\begin{equation}
\label{eq:conway_mononomials}
\psi^-_{i_1,0}\cdots\psi^-_{i_k,0}|s\rangle~,
\end{equation}
给出，这些单项式构成二十四维欧氏符号下的一个旋量。


下面我们想在上述态上构造 $\text{Co}_0$ 的作用。为此回忆 $\text{Co}_0\cong\Aut(\Lambda)$，故 Conway 群同构于 $\SO(24)$ 的子群，我们可以作自然等同 $\mathfrak{a}=\Lambda\otimes_\mathbb{Z}\mathbb{C}$，即让费米子“生活”在 Leech 格上。然后考虑本征值为 $\lambda_i^{\pm1}$ 的群元 $g\in\text{Co}_0$，选取 $\mathfrak{a}$ 的基使费米子 $\psi^\pm_i$ 按其本征向量变换：
\begin{equation}
\label{eq:conway_action_fermions}
g\psi^\pm_i=\lambda_i^{\pm1}\psi^\pm_i~,~~~\lambda_i\equiv e^{2\pi ia_i}~,~~~i=1,\dots,12~.
\end{equation}
此外，由于 Ramond 扇区的基态构成与 $\mathfrak{a}$ 相联系的 Clifford 代数的表示，我们应把 $G<SO(24)$ 提升为子群 $\widehat{G}<\text{Spin}(\mathfrak{a})$。元素 $x\in\text{Spin}(\mathfrak{a})$ 具有性质：对 $u\in\mathfrak{a}$ 有 $xux^{-1}\in\mathfrak{a}$。我们要求提升 $\hat{g}\in\widehat{G}<\text{Spin}(\mathfrak{a})$ 作用在 $\mathfrak{a}$ 上与 $g$ 相同，以此定义 $g\in G<\SO(\mathfrak{a})$ 的提升：
\begin{equation}
\label{lift_def}
\hat{g}(u):=\hat{g}u\hat{g}^{-1}=gu
~,~~~\forall~u\in\mathfrak{a}~.
\end{equation}
映射 $u\mapsto\hat{g}(u)$ 是 $\mathfrak{a}$ 上属于 $\SO(\mathfrak{a})$ 的线性变换，故 $\hat{g}\mapsto\hat{g}(\cdot)$ 定义了核为 $\lbrace\pm\mathbf{1}\rbrace$ 的映射 $\text{Spin}(\mathfrak{a})\rightarrow\SO(\mathfrak{a})$，即 $\text{Spin}(\mathfrak{a})$ 是 $\SO(\mathfrak{a})$ 的双重覆盖。事实证明，对 $G\cong\text{Co}_0$ 存在唯一的提升 $\widehat{G}\cong\Co_0$（细节见 \cite{DuncanJohnf.r.2015TMMF}）。

% [Tutor Note: 提升 $\hat{g}$ 的唯一性是 $\Co_0\subset\SO(24)$ 的特殊性质（对一般子群 $G<\SO(24)$ 并不成立）：$\text{Spin}(24)$ 的双重覆盖在 $\Co_0$ 上分裂，且分裂出的同构像在共轭意义下唯一。正因如此，下一段 $\nu$ 的符号（即 $\lambda_i^{1/2}$ 的选取）才有定义，否则 $T^s_g$ 会带整体符号歧义。]

NS 基态 $|0\rangle$ 在 $\Co_0$ 下不变，而群作用在 $2^{12}$ 个 Ramond 基态上则给出
\begin{equation}
\hat{g}|s\rangle=\prod_{i=1}^{12}e^{\pi i a_i}|s\rangle=\nu|s\rangle~,~~~\nu\equiv\prod_{i=1}^{12}\nu_i~,~~~\nu_i\equiv e^{\pi ia_i}=\lambda_i^{1/2}~,
\end{equation}
其中 $|s\rangle$ 是 \eqref{eq:conway_mononomials} 中描述的基态。注意先验地 $\nu_i$ 有符号歧义，因为它是 $\lambda_i$ 的平方根；但由于 $\text{Co}_0$ 的提升唯一，符号的选择实际上是唯一的。$g$ 在费米子上的作用还有另一种歧义：我们可以交换复本征值。这相当于在 \eqref{eq:conway_action_fermions} 中作替换 $\lambda_i\leftrightarrow-\lambda_i^{-1}$，称为极化（polarisation）的选取。


最后一步是考虑目前所描述理论的 $\ZZ_2=\lbrace1,\mathfrak{z}\rbrace$ 轨道流形，它在费米子上作用为 $\mathfrak{z}\psi^\pm_i=-\psi^\pm_i$。换言之，它作用为 $(-1)^F$，其中 $F$ 是费米子数。假定它在两个基态 $|0\rangle$ 与 $|s\rangle$ 上作用平凡\cite{DuncanJohnf.r.2015TMMF}，则两个扇区各自分解为偶/奇本征空间
\begin{equation}
\text{NS}=\text{NS}^0\oplus\text{NS}^1~,~~\text{R}=\text{R}^0\oplus\text{R}^1~,
\end{equation}
其中 $\text{NS}^j/\text{R}^j$ 的本征值为 $(-1)^j$。从此出发，可以构造两个密切相关的 SVOA。一个有用的描述利用了这一事实：$\text{NS}^0$ 自身构成一个（玻色）VOA，即格 $D_{12}$ 的 VOA。等价地，它是水平 $1$ 的仿射 Kac–Moody 代数 $\widehat{\mathfrak{so}}(24)_1$ 相联系的 VOA。后者有四个不可约可积模，即真空模 $A\cong\text{NS}^0$、向量模 $V$、旋量模 $S$ 与共轭旋量模 $C$。用任一旋量模扩张 $D_{12}$ VOA $\text{NS}^0$，就得到两个 SVOA：
\begin{equation}
\label{eq:Vf,Vs}
\begin{split}
&V^{f\natural}=\text{NS}^0\oplus S\cong\text{NS}^0\oplus \text{R}^0 \\
&V^{s\natural}=\text{NS}^0\oplus C\cong\text{NS}^0\oplus \text{R}^1~. 
\end{split}
\end{equation}
从轨道流形的观点看，在我们的记号下有等同 $S\cong\text{R}^0$ 与 $C\cong\text{R}^1$。


两个 SVOA $V^{f\natural}$ 与 $V^{s\natural}$ 作为 SVOA 同构，并由中心荷 $c_{V^{f\natural}}=c_{V^{s\natural}}=12$ 与无重量 $1/2$ 态唯一刻画。\cite{DuncanJohnF.2007SfCl} 证明，$V^{f\natural}$ 的 $\mathcal{N}=1$ 超流由 $\text{Spin}(24)$ 的、同构于 $\Co_0$ 的子群固定，该子群就是上述记号中的群 $\widehat{G}$。特别地，注意 $Z(\text{Spin}(24))\cong\ZZ_2\times\ZZ_2$，其中第一个 $\ZZ_2$ 因子可以等同为 $\text{Spin}(24) \to SO(24)$ 的核，后者等同为 $SO(24)$ 的中心。中心 $Z(\widehat{G})\cong\ZZ_2$ 可以等同为 Z(\text{Spin}(24)) 中的第二个 $\ZZ_2$，其作用与轨道流形构造中的 $\ZZ_2$ 相同。因此，由 \eqref{eq:Vf,Vs} 立即推出 $\Co_0$ 在 $V^{f\natural}$ 上的作用不是忠实的，因为后者在中心 $Z(\widehat{G})$ 的作用下不变。反之，$V^{f\natural}$ 带有商群 $\widehat{G}/\ZZ_2\cong\Co_1$ 的忠实作用。另一方面，$\Co_0$ 忠实地作用在 $V^{s\natural}$ 上，这最终正是我们后文取 $V^{s\natural}$ 而不取 $V^{f\natural}$ 的原因。$V^{f\natural}$ 与 $V^{s\natural}$ 之间另一个值得注意的差别是：$V^{f\natural}$ 中由 $\Co_0$ 固定的 $\mathcal{N}=1$ 超流并不包含在 $V^{s\natural}$ 中，而是包含在 $V^{s\natural}_{\text{tw}}$ 中（位于 $\text{R}^0$ 部分）。

% [Tutor Note: 尽管 $V^{f\natural}$ 与 $V^{s\natural}$ 作为 SVOA 同构，$\Co_0$ 在其上的作用却不同：$Z(\widehat{G})\cong\ZZ_2$ 在 $V^{f\natural}$ 上作用平凡，故只有商群 $\Co_1$ 忠实作用；而 $\Co_0$ 忠实作用在 $V^{s\natural}$ 上。这就是后文一律用 $V^{s\natural}$ 的原因——Conway 月光要的是 $\Co_0$（而非 $\Co_1$）的 McKay–Thompson 级数。]

$V^{s\natural}$ 的“正则扭模”（canonically twisted module，或 Ramond 扇区模）也可以构造为
\begin{equation}
V^{s\natural}_{\text{tw}}=\text{NS}^1\oplus\text{R}^0~,
\end{equation}
它是关于 $Z(\widehat{G})$ 对称性的扭模。$\Co_0$ 在这个扭模上的作用也是忠实的。


为了表述 Conway 月光，先定义函数
\begin{equation}
\label{eta_g_and_C_g}
\begin{split}
&\eta_{\pm g}(\tau):=q\prod_{n=1}^{\infty}\prod_{i=1}^{12}\left(1\mp\lambda_i^{-1}q^n\right)\left(1\mp\lambda_iq^n\right) \\
&C_{\pm g}:=\nu\prod_{i=1}^{12}\left(1\mp\lambda_i^{-1}\right)=\prod_{i=1}^{12}\left(\nu_i\mp\nu_i^{-1}\right)~.
\end{split}
\end{equation}
Conway 月光的 twined 配分函数则由下式给出
\begin{equation}
\begin{split}
T^s_g(\tau):=\text{str}_{V^{s\natural}}\left[\hat{g}q^{L_0-1/2}\right]=\text{tr}_{V^{s\natural}}\left[\mathfrak{z}\hat{g}q^{L_0-1/2}\right]=\frac{\eta_g(\tau/2)}{\eta_g(\tau)}+\chi_g \\
T^s_{g,\text{tw}}(\tau):=\text{str}_{V^{s\natural}_{\text{tw}}}\left[\hat{g}q^{L_0-1/2}\right]=\text{tr}_{V^{s\natural}_{\text{tw}}}\left[\mathfrak{z}\hat{g}q^{L_0-1/2}\right]=C_g\eta_g(\tau)-\chi_g~,
\end{split}
\end{equation}
其中超分次可以通过在迹中插入 $(-1)^F$ 来定义，其作用与 $\mathfrak{z}$ 相同（回忆我们把 $\ZZ_2=\lbrace1,\mathfrak{z}\rbrace$ 与 $\Co_0$ 的中心等同起来）。Conway 月光的主要定理指出 \cite{DuncanJohnf.r.2015TMMF}：

% [Tutor Note: $T^s_g$ 的上标 $s$ 是模的标记而非幂次：它表示迹取在 Conway 模 $V^{s\natural}$ 上，$s$ 来自模记号 $V^{s\natural}$，与 $V^{f\natural}$ 相对。$T^s_{g,\text{tw}}$ 则取在扭模 $V^{s\natural}_{\text{tw}}$ 上。两条迹实际都插入 $\mathfrak{z}\hat{g}$，故 $\text{str}$ 与 $\text{tr}$ 给出同一个量，不是两个独立的定义。]

% [Tutor Note: 式 \eqref{eq:eta_g_and_C_g} 中 $\eta_{\pm g}$ 的常数项：$\eta_g=q+O(q^2)$ 不含常数项，而 $\eta_g(\tau/2)/\eta_g(\tau)$ 的首项是 $q^{-1/2}$，其 $q^0$ 系数恰为 $-\chi_g$，故加上 $+\chi_g$ 后 $T^s_g$ 的常数项为零。这正是下一段定理中 $T^s_g(2\tau)=q^{-1}+O(q)$ 而非 $q^{-1/2}$ 的由来，因子 $2$ 不可漏掉。]

{\thm  \label{thm:conway_moonshine} 对每个 $g\in\Co_0$，函数 $T_g^s(2\tau)=q^{-1}+O(q)$ 是某个亏格零群 $\Gamma_g<\SL_2(\RR)$ 的主模，该群包含某个 $\Gamma_0(N)$。若 $g$ 在 $\Lambda$ 上的作用有不动点，则 $T^s_{g,\text{tw}}(\tau)$ 等于常数 $-\chi_g$。此外，若 $g$ 没有这样的不动点，则 $T^s_{g,\text{tw}}(\tau)$ 也是 $\SL_2(\mathbb{R})$ 的某个亏格零子群的主模。}\\

% [Tutor Note: 定理第二句的机制在定义式里就能看出：$C_g$ 应理解为 $C_{+g}=\prod_i(\nu_i-\nu_i^{-1})$；若 $g$ 有不动点，则某个 $\lambda_i=1$、$\nu_i=\pm1$，故 $C_g=0$，于是 $T^s_{g,\text{tw}}=C_g\eta_g(\tau)-\chi_g$ 退化为常数 $-\chi_g$。于是“扭模上的迹是否为常数”就成了“$g$ 在 $\Lambda$ 上是否有不动点”的判别。]

在 \S\ref{subsec:conwayK3} 中，我们将看到 Conway 月光在 K3 共形场论的对称性方面所起的突出作用。
